% The master problem (step 2), drawn as the incidence pattern of its rows:
% what a column writes into, what the artificials sit on, and -- the point --
% the row family that is not there.  The generation loop survives as the two
% margin arrows: columns enter top left, duals leave top right.
% Uses only tikz with calc/positioning/arrows.meta, all loaded by
% ../necessarypackages.  Drawn at final size (<= 15.1 cm against a 15.24 cm
% \textwidth): no \resizebox, so \scriptsize stays \scriptsize on the page.
% Every qualifier lives in the row label, never in a cell, so no cell can
% outgrow its row; the row heights below are set to the label line counts,
% and the column heads are at most three lines so they clear the head rule.
\begin{figure}[htb]
\centering
\hbadness=10000
\begin{tikzpicture}[x=1cm, y=1cm, >=Latex, line width=0.5pt,
  rl/.style  ={font=\scriptsize, align=right, text width=52mm, anchor=east},
  ce/.style  ={font=\small,      align=center},
  hd/.style  ={font=\footnotesize, align=center, text width=14mm,
               anchor=north, inner sep=1pt},
  nt/.style  ={font=\scriptsize, align=center, text width=18mm},
  dot/.style ={font=\small, color=black!45}]

% ---------------------------------------------------------------- geometry
\def\xS{5.55}\def\xB{12.85}                 % body left / right edge
\def\xa{6.28}\def\xb{7.74}\def\xc{9.20}\def\xd{10.66}\def\xe{12.12}
\def\xD{13.05}\def\xE{14.95}                % the dual strip
\def\yH{-2.25}
\def\yA{-2.95}\def\yB{-3.85}\def\yC{-4.75}\def\yD{-5.65}\def\yE{-6.80}
\def\yF{-7.70}\def\yG{-8.55}

% ------------------------------------------------------- shading and rules
\fill[black!7]  (\xS,\yH) rectangle (8.47,\yG);          % the column block
\fill[black!11] (\xD,\yA) rectangle (\xE,\yE);           % what pricing gets
\draw[line width=0.8pt] (0,\yH) -- (\xB,\yH);
\draw[line width=0.8pt] (0,\yG) -- (\xB,\yG);
\foreach \y in {\yA,\yB,\yC,\yD,\yE,\yF}
   \draw[black!35, line width=0.3pt] (0,\y) -- (\xB,\y);
\draw[black!35, line width=0.3pt] (\xS,\yH) -- (\xS,\yG);
\draw[black!70, line width=0.6pt] (8.47,\yH) -- (8.47,\yG);

% ------------------------------------------------------ the two pictograms
% a plant-bound load: stacks of full RTIs; a supplier-bound load: the same
% floor carrying folded ones.  That is the whole difference between the two
% directions as far as the master is concerned.
\draw[black!75, fill=white, line width=0.6pt]
   (\xa-0.66,-0.95) rectangle (\xa+0.66,-0.32);
\foreach \i in {0,1,2,3}{\foreach \j in {0,1}{
   \draw[black!70, fill=black!45, line width=0.3pt]
      (\xa-0.60+\i*0.31,-0.91+\j*0.255) rectangle ++(0.26,0.24);}}
\draw[black!75, fill=white, line width=0.6pt]
   (\xb-0.66,-0.95) rectangle (\xb+0.66,-0.32);
\foreach \i in {0,1,2,3}{\foreach \j in {0,1,2,3,4}{
   \draw[black!60, fill=black!18, line width=0.3pt]
      (\xb-0.60+\i*0.31,-0.91+\j*0.108) rectangle ++(0.26,0.09);}}

% --------------------------------------------- columns in, duals out (loop)
\node[rl, text width=46mm] at (4.78,-0.62)
   {each column comes from the pricing problem
    (Figure~\ref{fig:truck}): one truck load, with its layout};
\draw[->, line width=0.6pt] (4.90,-0.62) -- (5.50,-0.62);

\draw[->, line width=0.6pt] (14.00,\yA) -- (14.00,-1.74);
\node[nt, text width=21mm, anchor=south] at (13.95,-1.70)
   {to the pricing problem, rounded up};

% ------------------------------------------------------------ column heads
\node[hd] at (\xa,-1.10) {$\lambda^{p}_{\omega t}$\\[1pt]
   {\scriptsize full RTIs}};
\node[hd] at (\xb,-1.10) {$\lambda^{s}_{\omega' t}$\\[1pt]
   {\scriptsize folded RTIs}};
\node[hd] at (\xc,-1.10) {$\overline{V}_{nt}$,\\ $\overline{V^{\mathrm{e}}}_{rt}$\\[1pt]
   {\scriptsize aggregate}};
\node[hd] at (\xd,-1.10) {$I_{gnt}$,\\ $Ie_{gtr}$\\[1pt]
   {\scriptsize stock}};
\node[hd] at (\xe,-1.10) {artificial\\ columns$^{\dagger}$\\[1pt]
   {\scriptsize priced $M$}};

% ------------------------------------------------------------- row 1: cost
\node[rl] at (5.40,-2.60) {\textbf{objective} \eqref{eq:MPObjective}};
\node[ce]  at (\xa,-2.60) {$cK$};
\node[ce]  at (\xb,-2.60) {$cK$};
\node[dot] at (\xc,-2.60) {$\cdot$};
\node[ce]  at (\xd,-2.60) {$ic_{gn}$};
\node[ce]  at (\xe,-2.60) {$M$};

% ------------------------------------------ rows 2-3: the only link to load
\node[rl] at (5.40,-3.40) {\textbf{aggregate flow} \eqref{bp:MP-def-V}\\
   one row per product $n$, period $t$};
\node[ce]  at (\xa,-3.40) {$b_{\omega n}$};
\node[dot] at (\xb,-3.40) {$\cdot$};
\node[ce]  at (\xc,-3.40) {$-1$};
\node[dot] at (\xd,-3.40) {$\cdot$};
\node[ce]  at (\xe,-3.40) {$+1$};
\node[ce]  at (14.00,-3.40) {$\pi_{nt}$};

\node[rl] at (5.40,-4.30) {\textbf{aggregate flow} \eqref{bp:MP-def-Ve}\\
   one row per RTI type $r$, period $t$};
\node[dot] at (\xa,-4.30) {$\cdot$};
\node[ce]  at (\xb,-4.30) {$e_{\omega' r}$};
\node[ce]  at (\xc,-4.30) {$-1$};
\node[dot] at (\xd,-4.30) {$\cdot$};
\node[ce]  at (\xe,-4.30) {$+1$};
\node[ce]  at (14.00,-4.30) {$\pi_{rt}$};

% ------------------------------------------------------- row 4: truck count
\node[rl] at (5.40,-5.20) {\textbf{truck count} \eqref{bp:MP-count}\\
   one row per direction and period};
\node[ce]  at (\xa,-5.20) {$1$};
\node[ce]  at (\xb,-5.20) {$1$};
\node[dot] at (\xc,-5.20) {$\cdot$};
\node[dot] at (\xd,-5.20) {$\cdot$};
\node[dot] at (\xe,-5.20) {$\cdot$};
\node[ce]  at (14.00,-5.20) {$\mu_{gt}$};

% -------------------------------------------------------- row 5: cover cuts
\node[rl] at (5.40,-6.225) {\textbf{cumulative cover cut}
   \eqref{eq:covercut}$^{\dagger}$\\
   one row per direction and period, over the dispatches with $\tau\le t-lt$};
\node[ce]  at (\xa,-6.225) {$1$};
\node[ce]  at (\xb,-6.225) {$1$};
\node[dot] at (\xc,-6.225) {$\cdot$};
\node[dot] at (\xd,-6.225) {$\cdot$};
\node[ce]  at (\xe,-6.225) {$+1$};
\node[nt]  at (14.00,-6.225) {folds into $\mu_{gt}$};

% -------------------------------------------- rows 6-7: inventory, coverage
\node[rl] at (5.40,-7.25) {\textbf{inventory balance}
   \eqref{sb:eq:inv-plant-main}--\eqref{sb:eq:inve-sup-init}\\
   the receiver's arrival lagged by $lt$};
\node[dot] at (\xa,-7.25) {$\cdot$};
\node[dot] at (\xb,-7.25) {$\cdot$};
\node[ce]  at (\xc,-7.25) {$\pm1$};
\node[ce]  at (\xd,-7.25) {$\pm1$};
\node[dot] at (\xe,-7.25) {$\cdot$};

\node[rl] at (5.40,-8.125) {\textbf{coverage stock}
   \eqref{sb:eq:safety-full}--\eqref{sb:eq:safety-empty}};
\node[dot] at (\xa,-8.125) {$\cdot$};
\node[dot] at (\xb,-8.125) {$\cdot$};
\node[dot] at (\xc,-8.125) {$\cdot$};
\node[ce]  at (\xd,-8.125) {$+1$};
\node[ce]  at (\xe,-8.125) {$+1$};

\draw[black!55, line width=0.4pt] (12.95,\yE) -- (12.95,\yG);
\node[nt, text width=18mm, align=left, anchor=west] at (13.10,-7.675)
   {no $\lambda$: pricing feels these rows only through $\pi$};

% ------------------------------------------------- the rows that are absent
\node[rl] at (5.40,-9.45) {\textbf{rows a catalogue of loading\\
   configurations would add}};
\draw[dashed, black!55, line width=0.5pt] (\xS,-8.80) rectangle (\xB,-10.10);
\node[font=\scriptsize, align=center, text width=68mm] at (9.20,-9.45)
   {\textbf{none.} No slot, stack type, orientation, \mbox{coordinate} or
    layer count is a variable of the master: the geometry stays inside the
    \mbox{column}, with its witnessing layout (Figure~\ref{fig:truck}).};

% ----------------------------------------------------------------- the note
\node[font=\scriptsize, align=left, text width=149mm, anchor=north west]
   at (0,-10.38)
   {$^{\dagger}$ Rows and columns of the \emph{restricted} master;
    \eqref{eq:MPObjective}--\eqref{bp:MP-count} together with
    \eqref{sb:eq:inv-plant-main}--\eqref{sb:eq:safety-empty} are the master
    itself.};

\end{tikzpicture}
\caption{The master is a flow model that has never seen a box. A column is
one truck load, and the only rows it enters are the two aggregate-flow
definitions, one truck count per direction and period, and the cumulative
cover cut --- every one of which merely counts or sums what the load
contains. No slot, stack type, orientation, coordinate or layer count is a
variable here, so there is no catalogue of loading configurations to
enumerate, and no row written over these variables can reach the geometry;
the inventory and coverage rows see the trucks only through $\overline{V}$
and $\overline{V^{\mathrm{e}}}$, and all that ever reaches the pricing
problem is $\pi$ and $\mu$. The non-negative artificials keep the restricted
master feasible, and its duals defined, for any pool --- a solution that uses
them is not a plan.}
\label{fig:loop}
\end{figure}
